\section{Sugestões para Trabalhos Futuros}
\begin{itemize}
    \item Implementação de uma bancada física com 6 graus de liberdade;
    \item Caracterização da bateria (curva SOC $\times$ OCV);
    \item Extensão do observador para a região de enfraquecimento de campo,
    fora do envelope de operação ensaiado neste trabalho;
    \item Uso de dinamômetro estático ou dinâmico para medição direta de torque
    no eixo do motor, dispensando a estimativa do observador nos ensaios de
    perdas mecânicas e permitindo separar as parcelas de atrito seco e histerese
    que a regressão linear atual agrega em $\tau_c$;
    \item Estudo do impacto da temperatura da bateria na tração do conjunto
    motopropulsor, relacionando a queda de tensão e a resistência interna do
    pacote ao longo do ensaio com a rotação e o empuxo medidos.
    \item Repetição dos ensaios das hélices com o levantamento das perdas a
    vazio no mesmo dia e com o mesmo controlador, e com o registro da pressão e
    da temperatura do ambiente, para separar as parcelas da diferença em relação
    ao fabricante;
    \item Envio ao computador de lotes inteiros de amostras do conversor, que já
    passam pela memória a \SI{111}{\kilo\hertz}, para obter as formas de onda de
    tensão e corrente pela própria placa, sem osciloscópio;
    \item Configuração automática da rede, com a placa atribuindo o endereço do
    computador, para dispensar ajustes de endereço e de \textit{firewall} a cada
    conexão;
    \item Ensaios em túnel de vento, para caracterizar as hélices fora da
    condição estática.
\end{itemize}
